% ============================================================
% ملخص — Arabic Abstract
% Requires XeLaTeX + polyglossia + an Arabic-capable font (Amiri or Scheherazade)
% TODO: If compilation fails due to missing Arabic font, install Amiri:
%       https://github.com/alif-type/amiri/releases
%       or change \arabicfont to a system-installed Arabic font.
% ============================================================

\chapter*{ملخص}
\addcontentsline{toc}{chapter}{ملخص}

\begin{Arabic}

يُعدّ الرصد الزراعي في البيئات الصحراوية تحدياً متعدد الأوجه يختلف جوهرياً عن سياقات
الزراعة التقليدية؛ إذ تتسم هذه البيئات بانقطاع التيار الكهربائي، وعدم استقرار
الاتصال بالإنترنت، وقسوة الظروف الحرارية، فضلاً عن صعوبة الوصول المنتظم إلى
القطع الزراعية المتفرقة. وفيما يخص زراعة الفصة تحديداً، تتفاقم هذه القيود بسبب
شُح أنظمة مراقبة إنترنت الأشياء المُصمَّمة خصيصاً لهذا المحصول في السياق الصحراوي
الجزائري، إضافةً إلى قصور متكرر تم رصده في الأعمال العلمية المُراجَعة، يتمثل في
الخلط بين توقيت رفع البيانات وتوقيت قياسها الفعلي، مما يُضعف موثوقية أي تحليل
زمني مبني على هذه البيانات.

تقدم هذه الرسالة تصميم وتنفيذ منصة متكاملة لرصد الزراعة الذكية وتحليلها الحتمي،
مُطبَّقة على زراعة الفصة في منطقة الوادي بالجزائر. يتكون النظام من ثلاثة عقد
مدمجة تتواصل عبر ارتباط لاسلكي LoRa على تردد 433~ميغاهرتز وفق إطار زمني دوري
مدته عشر دقائق. تتولى العقدة البوابة الرئيسية تنسيق استقبال الحزم، والحفاظ على
توقيت مرجعي دقيق عبر ساعة DS3231 الزمنية، وتخزين جميع القياسات محلياً على بطاقة
ذاكرة SD بمعزل عن توفر الشبكة. تجمع العقدة الثانية (Node2) بيانات رطوبة التربة
ودرجة حرارتها وموصليتها الكهربائية عند محور الري الثاني عبر بروتوكول RS485/Modbus،
بينما ترصد العقدة الثالثة (Node3) درجة الحرارة والرطوبة النسبية والضغط الجوي.
يعتمد نقل البيانات مبدأ الأولوية المحلية (Offline-First)، إذ تُراكَم السجلات
على البطاقة الذاكرة وتُرسَل دُفعةً واحدة إلى الخادم البعيد حين يتوفر اتصال الإنترنت.

تُعالج ثمانية أوضاع تشغيلية للبرنامج الثابت قيود الميدان المحددة، من بينها:
آلية استرداد تُصحح الانجراف الزمني المسؤول عن أعطال الاتصال المتكررة، وأوضاع
توفير الطاقة التي تُطيل عمر البطاريات عبر النوم العميق والتحكم في تغذية
الحساسات بالمرحّلات. يُطبّق الخادم الخلفي المبني على Node.js وقاعدة بيانات
PostgreSQL سياسةً صارمة تفصل الختم الزمني المستمَد من ساعة الوقت الحقيقي
(measured\_at) عن توقيت استقبال الرفع، المحفوظ حصراً بوصفه
سجلاً تدقيقياً. يعرض لوح التحكم المبني بإطار React اتجاهات الحساسات والأحداث
الزراعية ونتائج خط أنابيب تحليلي حتمي من عشر مراحل يشمل ضبط الجودة وتقييم
الموثوقية. يعمل نظام التفسير المستند إلى الذكاء الاصطناعي (Gemini) عبر وسيط
خلفي خاضع لستة قيود معمارية تحول دون قيامه بحسابات التحليل الأولية أو تعديل
البيانات المخزنة.

أثبتت نافذة التحقق التجريبي الممتدة من 6 مايو 2026 إلى 11 مايو 2026 نجاح
المنصة، إذ سُجِّل 2433 قراءة استشعارية موزّعة بالتساوي على العقد الثلاث، وأُنجزت
ست جلسات رفع بالكامل، وأُنتج 135 صفاً من مخرجات المعالجة الحتمية. تُثبت هذه
النتائج الجدوى الوظيفية للمعمارية الشاملة في إطار نشر تجريبي أولي؛ في حين تُعد
التحقق الموسمي الكامل ومعايرة عتبات الموثوقية أولوياتٍ للأعمال البحثية المستقبلية.

\vspace{0.8em}
\noindent\textbf{الكلمات المفتاحية:}
إنترنت الأشياء، الزراعة الذكية، الفصة، LoRa، ESP32، الأولوية المحلية،
التحليل الحتمي، الزراعة الصحراوية، جودة البيانات، ساعة الوقت الحقيقي،
تتبع التوقيت الزمني.

\end{Arabic}
